\subsection{Checks for \texorpdfstring{Section~\ref{sec:ren}}{renormalization}}\label{chk:ren}
% src: H4_rg_static, H4_hier, H4_memory, H4_codes, H4_lrg outputs; independent recomputations C3_04_rg, C3_05_memory, C3_09_latent; H1_checks (check A); H6_nested, H6_misc (nested media under decision 22); ren_checks.py (parts A to D) written while drafting this section; pvrf_ren.py (prop:ren-per-level(f), verification round); ren_nonuniform_leak.py (thm:ren-memory(d) and thm:ren-closure(a) with nonuniform leaks, 2026-11-10)

Each check below evaluates exact Gaussian linear algebra (Schur complements, Lyapunov solutions, Fourier inversion on tori, spectral decompositions) on the instances stated. Where a second, independent implementation of the same statement was run, its results are reported in the same check.

\begin{numcheck}[Decimation on $\Z$]\label{chk:ren-decimation}
Rings of $60b$ sites with $(\varepsilon^\Pi,a^\Pi)\in\{(0.01,1),(0.3,0.7),(2,0.5)\}$ and $b\in\{2,3,5\}$ were decimated by Schur complements. The coarse couplings and leaks agreed with \eqref{eq:ren-decimation} to $13$ digits, for instance $a^{\Pi,(1)}=0.4975124378$ and $\varepsilon^{\Pi,(1)}=0.01995024876$ at $(0.01,1)$ and $b=2$; every coupling beyond range one vanished exactly, and $m^{(1)}/m=b$ to $12$ digits. The map $\mathfrak x\mapsto\mathfrak x(\mathfrak x+4)$ reproduced $\mathfrak x^{(1)}=0.0401$ at $\mathfrak x=0.01$. The squared second-moment length of the decimated field agreed with $1/(4\sinh^2(bm/2))$ to $10^{-8}$ relative at $\varepsilon^\Pi\in\{10^{-3},10^{-2},0.1\}$ and $b\in\{2,4\}$. An independent implementation found the largest relative deviation $3.2\times10^{-13}$ over the three media.
\end{numcheck}

\begin{numcheck}[Nearest-neighbor passing on $\Z$]\label{chk:ren-ring}
For $(c,a,q)=(0.25,0.3,1)$ and $(0.1,0.44,0.7)$ and $b\in\{2,3,4\}$ on a ring of $400$ sites, the two-pole formula of \cref{prop:ren-ring}(a) matched the decimated covariance to $1.7\times10^{-14}$ relative, the largest off-diagonal entry of every decimated precision was negative, and the ratios of consecutive tail couplings, $0.120538$, $0.024868$, $0.021248$ for the first medium and $0.363636$, $0.053705$, $0.152641$ for the second, equaled $e^{-\arccosh u_\ast}$ to six digits, the odd factor $b=3$ included.
\end{numcheck}

\begin{numcheck}[Decimation on $\Z^2$]\label{chk:ren-two-dim}
On a $16\times16$ torus with $\varepsilon^\Pi=0.05$ and $a^\Pi=1$ the checkerboard decimation gave the couplings $0.493827$ and $0.246914$ at the rotated nearest and diagonal displacements and the leak $0.0993827$, as \cref{ex:ren-two-dim} predicts. The further decimation onto $2\Z\times2\Z$ gave all $63$ couplings per site positive, decaying from $0.4575$ at the coarse displacement $(1,0)$ to $0.1154$ at $(1,1)$, $0.0189$ at $(2,0)$ and $0.0017$ at $(3,0)$. For \cref{obs:ren-decay} the covariance of the nearest-neighbor field with $a^\Pi=1$ on a $256\times256$ torus was computed from its symbol by Fourier inversion, subsampled on $2\Z\times2\Z$ and on $4\Z\times4\Z$, and inverted by Fourier transform. At $\varepsilon^\Pi=0.5$, $0.05$, $0.005$ and $0.0005$ the coarse couplings at $(1,0)$ to $(5,0)$ after the first decimation were $(0.3565,8.92\times10^{-3},5.95\times10^{-4},3.98\times10^{-5},2.67\times10^{-6})$, $(0.4575,1.89\times10^{-2},1.71\times10^{-3},1.55\times10^{-4},1.41\times10^{-5})$, $(0.4712,2.06\times10^{-2},1.92\times10^{-3},1.81\times10^{-4},1.71\times10^{-5})$ and $(0.4726,2.07\times10^{-2},1.95\times10^{-3},1.84\times10^{-4},1.74\times10^{-5})$, with consecutive ratios between $0.0250$ and $0.0948$; after the second they were $(0.0698,1.34\times10^{-3},3.71\times10^{-5},1.08\times10^{-6},3.21\times10^{-8})$, $(0.2349,1.87\times10^{-2},1.85\times10^{-3},1.91\times10^{-4},2.01\times10^{-5})$, $(0.2844,2.76\times10^{-2},3.28\times10^{-3},4.06\times10^{-4},5.13\times10^{-5})$ and $(0.2903,2.88\times10^{-2},3.49\times10^{-3},4.39\times10^{-4},5.66\times10^{-5})$, with ratios between $0.0192$ and $0.1288$. The ratios increase with the displacement: at $\varepsilon^\Pi=0.0005$ the ratios of the couplings at $(x+1,0)$ and $(x,0)$ for $x=1,\dots,8$ were $0.0439$, $0.0938$, $0.0943$, $0.0948$, $0.0953$, $0.0957$, $0.0961$, $0.0964$ after the first decimation and $0.0993$, $0.1210$, $0.1260$, $0.1288$, $0.1309$, $0.1324$, $0.1336$, $0.1346$ after the second, the same to the printed digits on a $512\times512$ torus. No coarse precision had a positive off-diagonal entry.
\end{numcheck}

\begin{numcheck}[Decay of decimated kernels]\label{chk:ren-decay-rate}
The decimated precision of the nearest-neighbor precision, with $a^\Pi=1$ except where another value is stated, was computed on tori of $32^2$, $30^2$ or $40^3$ coarse sites from the symbol $1/\hat\Sigma^{\mathsf{Dec}}$ of \eqref{eq:ren-symbols}, which gives the periodization of the couplings on the infinite lattice, and compared with \cref{prop:ren-decay-rate} (script \texttt{check\_c.py}, deterministic). For $b=2$ on $\Z^2$ at $\varepsilon^\Pi=0.5$, $0.05$, $0.005$ and $0.0005$, $e^{-\mu_1}=0.070524$, $0.097036$, $0.100610$ and $0.100979$ from \eqref{eq:ren-decay-two}, with the limit $0.101021$ at $\varepsilon^\Pi=0$; the ratio of the hyperplane sums at coarse distances six and five agreed with these values to six digits, and the hyperplane sums agreed with \eqref{eq:ren-decay-hyp} and with the closed form of \eqref{eq:ren-decay-two} to $8.5\times10^{-13}$ relative at distances two to four and to $4.2\times10^{-9}$ at distances up to seven. The loss at the larger distances is the rounding of the transform: there the sums are as small as $1.4\times10^{-8}$ times the sum at distance zero, and the absolute deviation at distances two to seven was at most $1.5\times10^{-16}$ times that sum. For $b=4$ on $\Z^2$ at the same leaks, $e^{-\mu_1}=0.032311$, $0.116244$, $0.139727$ and $0.142477$, with $\mu_1=3.432$, $2.152$, $1.968$ and $1.949$ between $bm_1=2.773$, $0.893$, $0.283$ and $0.089$ and $bm_\ast=5.802$, $5.325$, $5.274$ and $5.268$, and agreement of the hyperplane sums with \eqref{eq:ren-decay-hyp} to $1.9\times10^{-11}$ relative at distances two to four and to $5.2\times10^{-8}$ at distances up to seven, where the sums fall to $3.2\times10^{-11}$ times the sum at distance zero. At $\varepsilon^\Pi=0.05$, $e^{-\mu_1}=0.108984$ for $b=3$ on $\Z^2$, $0.169894$ for $b=2$ on $\Z^3$ and $0.222133$ for $b=3$ on $\Z^3$, with agreement to $1.3\times10^{-10}$ or better at distances two to seven. Over the eleven cases the absolute deviation at distances two to seven was at most $2.9\times10^{-16}$ times the sum at distance zero, and the value of the hyperplane sum at distance one given in \cref{prop:ren-decay-rate}(a) held to $2.1\times10^{-15}$ relative. In every case the largest off-diagonal entry of the decimated precision was below $10^{-16}$, that is nonpositive up to rounding, and the couplings on the axis at distances one to seven were at most the hyperplane sums in absolute value. The ratios of successive couplings on the axis lay below $e^{-\mu_1}$ and rose with distance: for $b=2$ on $\Z^2$ at $\varepsilon^\Pi=0.0005$ they were $0.0439$, $0.0938$, $0.0943$, $0.0948$ and $0.0953$ from distances one and two to distances five and six, and for $b=4$ they were $0.0993$, $0.1210$, $0.1260$, $0.1288$ and $0.1309$, in agreement with the factors recorded in \cref{obs:ren-decay}. The last of these ratios exceeded $e^{-\mu_1}(5/6)^{(D-1)/2}$ by between $0.2$ and $4.5$ percent over the eleven cases, which is consistent with a correction $n^{-(D-1)/2}$ to the exponential and does not prove it. For part (d), the partial transforms of the couplings over $x'_\perp$ at all transverse momenta of the torus and at $n=2$, $3$ and $4$ were compared with \eqref{eq:ren-decay-partial} for $(b,D,\varepsilon^\Pi)=(2,2,0.05)$, $(4,2,0.05)$, $(3,2,0.5)$ and $(2,3,0.05)$, on tori of $32^2$, $32^2$, $30^2$ and $24^3$ coarse sites: they agreed to $7.6\times10^{-16}$ times the hyperplane sum at distance zero, and to $1.2\times10^{-10}$ relative where they exceeded $10^{-6}$ times that sum; none was positive beyond $2.2\times10^{-17}$ times that sum; no rate $\mu_i(p_\perp)$ on these grids was below $\mu_1$; and $\mathsf H(\cdot;p_\perp)$ had a single pole at one momentum in each of the two cases with $b=2$ and at none in the others. The couplings on the axis were then computed at large distances from \eqref{eq:ren-decay-axis-integral}, by the rectangle rule on $4096$ transverse momenta on $\Z^2$ and $160^2$ on $\Z^3$ and in logarithms; at distances two to four they agreed with the torus values to $8.7\times10^{-13}$ relative. For $b=2$ on $\Z^2$ at $\varepsilon^\Pi=0.0005$ the local rate $-\ln(\Pi^{\mathsf{Dec}}((n+1)\mathbf e_1)/\Pi^{\mathsf{Dec}}(n\mathbf e_1))$ minus $\mu_1$ was $0.039989$, $0.023092$, $0.012092$, $0.006154$, $0.003102$ and $0.001557$ at $n=10$, $20$, $40$, $80$, $160$ and $320$, against $\frac12\ln(1+1/n)=0.047655$, $0.024395$, $0.012346$, $0.006211$, $0.003115$ and $0.001560$; at $n=320$ it was $0.001559$ for $b=4$ at the same leak and $0.001554$ for $b=2$ at $\varepsilon^\Pi=0.5$; and for $b=2$ on $\Z^3$ at $\varepsilon^\Pi=0.05$ it was $0.094476$, $0.048726$ and $0.024690$ at $n=10$, $20$ and $40$, against $\ln(1+1/n)=0.095310$, $0.048790$ and $0.024693$. The quantity $-n^{-1}\ln(-\Pi^{\mathsf{Dec}}(n\mathbf e_1))-\mu_1$, which tends to zero by \eqref{eq:ren-decay-axis}, was $0.0044$, $0.0065$ and $0.0044$ at $n=320$ in the three cases on $\Z^2$, having decreased from $n=40$ on, and $0.186$, $0.127$ and $0.081$ at $n=10$, $20$ and $40$ on $\Z^3$. The agreement of the local rates with $\frac{D-1}2\ln(1+1/n)$ supports the correction $n^{-(D-1)/2}$ and does not prove it; for $b=2$ on $\Z^2$ the correction is proved in part (e). The dependence on the coupling was tested with $(b,D,\varepsilon^\Pi,a^\Pi)=(2,2,0.3,0.7)$, $(4,2,0.02,2.5)$, $(3,2,0.2,1.9)$ and $(2,3,0.1,1.3)$ on tori of $32^2$, $32^2$, $30^2$ and $40^3$ coarse sites. There $e^{-\mu_1}=0.073864$, $0.137940$, $0.097374$ and $0.165093$, the ratio of the hyperplane sums at distances six and five agreed with these values to five digits, the hyperplane sums agreed with \eqref{eq:ren-decay-hyp} to $5.7\times10^{-13}$ relative at distances two to four and to $1.8\times10^{-9}$ at distances up to seven, with an absolute deviation of at most $1.3\times10^{-16}$ times the sum at distance zero, the value at distance one held to $6.6\times10^{-16}$ relative, the largest off-diagonal entry was below $10^{-16}$, the axis couplings at distances one to seven were at most the hyperplane sums in absolute value, and the difference between the two sides of \eqref{eq:ren-decay-gap} was $0$, $1.40$, $0$ and $0.561$, the two zeros being the cases with exactly two poles, where the bound is an equality. The closed form \eqref{eq:ren-decay-two} with $(\varepsilon^\Pi,a^\Pi)=(0.3,0.7)$, $(0.02,2.5)$ and $(1.1,0.4)$ held to $2.1\times10^{-11}$ relative at distances two to four, and the partial transforms at all transverse momenta of the torus and $n=2$, $3$ and $4$ agreed with \eqref{eq:ren-decay-partial} to $4.2\times10^{-16}$ times the hyperplane sum at distance zero for the first, third and fourth of the four cases, the last on a torus of $24^3$ coarse sites. For part (e), with $b=2$ on $\Z^2$ and $(\varepsilon^\Pi,a^\Pi)=(0.0005,1)$, $(0.5,1)$, $(0.3,0.7)$ and $(0.02,2.5)$, the symbol of the decimated precision on the torus of $32^2$ coarse sites agreed with \eqref{eq:ren-decay-two-symbol} to $6.6\times10^{-14}$ relative, the partial transforms at all $32$ transverse momenta and $n=2$, $3$ and $4$ agreed with \eqref{eq:ren-decay-two-partial} to $9.0\times10^{-17}$ times the hyperplane sum at distance zero and to $1.4\times10^{-11}$ relative where they exceeded $10^{-6}$ times that sum, and the axis couplings obtained by integrating \eqref{eq:ren-decay-two-partial} over $4096$ momenta agreed with the torus values at distances two to four to $1.4\times10^{-12}$ relative. The limits of \eqref{eq:ren-decay-two-axis} are $4.325821$, $4.339189$, $3.035424$ and $10.814566$ in the four cases. In the first, $n^{1/2}e^{\mu_1n}(-\Pi^{\mathsf{Dec}}(n\mathbf e_1))$ was $4.084626$, $4.234536$, $4.285756$, $4.306890$, $4.316607$, $4.321275$ and $4.323562$ at $n=10$, $20$, $40$, $80$, $160$, $320$ and $640$, and $n$ times its relative distance from the limit was $0.5576$, $0.4220$, $0.3705$, $0.3501$, $0.3408$, $0.3363$ and $0.3341$, which levels off as a relative correction of order $1/n$ does. In the other three cases the same product was $0.5778$, $0.5403$ and $0.3374$ at $n=640$, after $0.5818$, $0.5440$ and $0.3396$ at $n=320$.
\end{numcheck}

\begin{numcheck}[Flows of the second-moment length and of the axis rate]\label{chk:ren-flows}
In one dimension, for the nearest-neighbor precision at $\varepsilon^\Pi\in\{10^{-3},10^{-2},0.1\}$ and $b\in\{2,4\}$, both formulas of \eqref{eq:ren-xi2-flow} agreed with the covariance-side lengths of the coarse fields to $10^{-8}$ relative; at $\varepsilon^\Pi=10^{-2}$ and $b=2$ the squared lengths were $24.93766$ (decimation) and $25.06234$ (averaging), against $\xi_2^2/b^2=25$. In two dimensions, on a $480\times480$ torus at $\varepsilon^\Pi\in\{0.01,0.1\}$ and $b\in\{2,3,4\}$, they agreed to all eight printed digits; at $\varepsilon^\Pi=0.01$ and $b=4$ the lengths were $2.44809$ and $2.51551$ against $\xi_2/b=2.5$, and at $\varepsilon^\Pi=0.1$ and $b=4$ they were $0.66218$ and $0.83652$ against $0.79057$. The decimated length lay below $\xi_2/b$ in all six two-dimensional cases. On a $510\times510$ torus, for passing with $(c,a)=(0.2,0.19)$ and $(0.3,0.17)$, the axis rates of the decimated and of the block-averaged fields equaled $bm_1$ to five digits or better for $b=2,3$, with relative deviation at most $1.6\times10^{-6}$; for instance $m_1=0.45489907$, and at $b=2$ the decimated field gave $0.90979956$ and the averaged one $0.90979814$, against $2m_1=0.90979814$. An independent implementation reproduced the one-dimensional formulas to $1.7\times10^{-9}$ on a ring of $24000$ sites and the two-dimensional values to every printed digit.
\end{numcheck}

\begin{numcheck}[Block averages]\label{chk:ren-block-average}
On a ring of $1200$ sites with $\varepsilon^\Pi\in\{10^{-3},10^{-2},0.2\}$ and $b\in\{2,3,4,8\}$, the ratio of the precision entries of the block averages at distances two and one equaled $-\mathsf h_b$ of \cref{ex:ren-block-average} to six digits in all twelve cases, with a positive entry at distance two; at $\varepsilon^\Pi=10^{-3}$ and $b=2$ it was $-0.171512$, against the limit $-0.171573$. An independent implementation found the largest relative error $3.8\times10^{-14}$ and the limits $0.1716$ and $2-\sqrt3$. On a $240\times240$ torus with the nearest-neighbor kernel, $b\in\{2,3\}$ and $\varepsilon^\Pi\in\{10^{-3},0.05\}$, the block-average precision had positive off-diagonal entries: counted per site, as the number of nonzero coarse displacements at which the entry of the precision kernel exceeds $10^{-12}$, there were $316$ and $308$ for $b=2$ and $376$ and $372$ for $b=3$, at the smaller and the larger leak, and the same four counts on a $480\times480$ torus. The entries decay with the displacement, so the count grows as the threshold is lowered, from between $144$ and $184$ at the threshold $10^{-8}$ to between $216$ and $268$ at $10^{-10}$; below about $10^{-14}$ rounding dominates. At $\varepsilon^\Pi=10^{-3}$ its coupling at the coarse displacement $(2,0)$ was $-0.272$ for $b=2$ and $-0.451$ for $b=3$, against $2.031$ and $2.590$ at $(1,0)$.
\end{numcheck}

\begin{numcheck}[Long-range couplings]\label{chk:ren-exponents}
The kernel $a^\Pi(x)=1/x^2$ on $\Z$ has the symbol $\varepsilon^\Pi+\pi\lvert p\rvert-p^2/2$ on $[-\pi,\pi]$, so $\sigma=1$, $\sigma'=2$ and $\mathsf a_\sigma=\pi$. The coarse symbols \eqref{eq:ren-symbols} were evaluated in closed form at $p'=0$, $h$ and $2h$ with $h=10^{-5}\varepsilon^\Pi$, and $\varepsilon^{\Pi,(1)}$ and $\mathsf a^{(1)}_\sigma$ were extracted from the first two orders. At $\varepsilon^\Pi\in\{10^{-2},10^{-3}\}$ and $b\in\{2,3,5\}$, $\varepsilon^{\Pi,(1)}$ agreed with \cref{prop:ren-exponents}(b) to ten digits for both maps and $\mathsf a^{(1)}_\sigma$ to $2\times10^{-7}$ relative; the factor by which $\varepsilon^\Pi/\mathsf a_\sigma$ was multiplied was $2.00404$, $3.01365$ and $5.05263$ for decimation at $\varepsilon^\Pi=10^{-2}$, against $b(1+\varepsilon^\Pi\mathsf S_1)$, and $2.00000$, $3.00000$ and $5.00000$ for averaging.
\end{numcheck}

\begin{numcheck}[The nested medium]\label{chk:ren-nested}
For $(b,n_{\mathrm H},A_{\mathrm h},\mathfrak q,\varepsilon^\Pi)=(3,5,1,0.5,10^{-3})$, $(2,8,1,0.7,10^{-4})$ and $(4,4,2,0.3,0.05)$, with $243$, $256$ and $256$ agents, $\Vert\Pi_{\perp\Phi}\Vert_F\le1.6\times10^{-14}$, the coarse precision equaled the nested precision with amplitude $\mathfrak qA_{\mathrm h}$ to $1.4\times10^{-15}$, the inverse covariance of the block coordinates equaled $\Pi_\Phi$ to $1.1\times10^{-12}$, and the contrast of the blocks of every level agreed to ten digits in the fine medium, by \eqref{eq:ren-nested-contrast} and in the coarse medium; the third medium has $r=3.073152$, $2.992958$ and $2.788462$ for blocks of levels one to three. For this medium and $r_0=2$, the formula of \cref{thm:ren-nested}(d) gives $l^\ast=2$, and blocks of levels one to three reach $r_0$. In finite hierarchies the top levels have small cuts: in the second medium $r$ rises from $0.934$ at level one to $2.855$ at level seven. An independent implementation on the media $(3,4,1,0.5,10^{-3})$, $(2,7,1,0.7,10^{-4})$ and $(4,3,2,0.3,0.05)$ found the equitability residual at most $7.3\times10^{-16}$, the error of the flow at most $8.9\times10^{-16}$ and equal contrasts in the fine medium, by the formula and in the coarse medium.
\end{numcheck}

\begin{numcheck}[Depth]\label{chk:ren-depth}
For decimation of the nearest-neighbor precision on $\Z$ with $a^\Pi=1$ and $b=2$, the largest block lengths $2^l$ at which the necessary condition of \cref{prop:ren-depth}(a) holds, at $(r_0,\varepsilon^\Pi)=(2,10^{-4})$, $(2,10^{-3})$, $(5,10^{-4})$ and $(5,10^{-3})$, were $128$, $32$, $64$ and $16$, below the bounds $(2/m)\arcsinh((r_0-1)^{-1/2})=176.3$, $55.7$, $96.2$ and $30.4$ and $2\xi_{\mathrm{OZ}}/(r_0-1)^{1/2}=200.0$, $63.2$, $100.0$ and $31.6$, and at the next level the best adjacent pair had contrast $1.083$, $1.149$, $1.348$ and $1.473$, below $r_0$. The condition is only necessary: in these homogeneous media every region has contrast below $2$, an adjacent pair having $(\mathfrak x^{(l)}+2)/(\mathfrak x^{(l)}+1)$, so no set reaches either level. For \cref{prop:ren-depth}(d), $60$ random symmetric translation-invariant kernels on rings, of range one to six, with leaks in $[10^{-3},1]$ and blocks of two to five sites, gave coarse-medium lengths that agreed with the formula to $4.5\times10^{-13}$ relative and never exceeded $\xi_2^2/B$; nearest-neighbor kernels at $\varepsilon^\Pi=10^{-3}$ gave $500$, $333.33$, $250$ and $125$ for $B=2$, $3$, $4$ and $8$, equal to $\xi_2^2/B$. Six random kernels of range two on tori of $12\times12$ square blocks of side two or three agreed with the formula to $6.3\times10^{-14}$, with ratios to $\xi_2^2/B$ at most $0.715$.
\end{numcheck}

\begin{numcheck}[Nested collectives level by level]\label{chk:ren-per-level}
In the coarse media built by iterating $\mathsf Q\Pi^{(l)}\mathsf Q^\top$, with $\mathsf Q$ the indicator matrix of the next partition, for $(b,n_{\mathrm H},A_{\mathrm h},\mathfrak q,\varepsilon^\Pi)$ equal to $(3,5,1,0.5,10^{-3})$, $(2,8,1,0.7,10^{-4})$, $(4,4,2,0.3,0.05)$ and $(2,8,1,0.2,0.7)$, the contrast of the blocks of the next level and the exchangeable baseline agreed with the closed forms at every level, the baseline to $5\times10^{-7}$ at values up to $3.2\times10^4$. The first three media have $\mathfrak q\ge1/b$ and $\Delta_{\mathrm H}<0$ at every level, and no level qualified at any of $121$ leaks in $[10^{-4},10^2]$. In the fourth the first qualifying leak of the grid at each level $l=0,\dots,6$ was the first grid point above the threshold $A_{\mathrm h}\mathfrak q^lT_{\mathrm H}$, which is $0.6336$, $0.1285$, $0.02643$, $0.005593$, $0.001255$, $3.174\times10^{-4}$ and $1.024\times10^{-4}$; at $\varepsilon^\Pi=1$ the windows of admissible levels, $(2.2549,2.4000]$, $(1.2520,1.3333]$, $(1.0508,1.0693]$ and so on down to $(1.00010,1.00012]$, are disjoint, as (e) requires. The sign of $\Delta_{\mathrm H}$ on $260739$ triples ($b=2,\dots,10$, $999$ values of $\mathfrak q$, $m=2,\dots,30$) showed no exception to (c) or (d). For the embedding, the sixteen agents of the medium with $b=2$, $n_{\mathrm H}=4$ and $A_{\mathrm h}=1$ were placed at the sites $0,\dots,15$ of $\Z$, each level was given the lattice clause of \cref{def:col-baseline} on its own box, and every subset of every comparison box was enumerated. Judged level by level, the blocks of all three levels are collectives for $\varepsilon^\Pi$ above $0.00220$, $0.0107$, $0.0399$, $0.780$ and $3.44$ at $\mathfrak q=0.1$, $0.2$, $0.3$, $0.5$ and $0.55$; at $\mathfrak q=0.6$ the third level never qualifies, and at $\mathfrak q=0.7$ no level does.

For (f), $\Delta_{\mathrm H}(b,1/(2b-1),m)$ was computed in exact rational arithmetic for $b=2,\dots,12$ and $m=2,\dots,60$: every value is negative and equals $(2b-1)^{1-m}/(b-1)-S_m/(b^m-1)$, with $-1/((2b-1)(b+1))$ at $m=2$, for instance $-0.1111$ for $b=2$. The strict increase in $m$, with the numerator identity of the proof, held at all $8120$ exact triples with $b=2,\dots,8$, $\mathfrak q=j/(40(2b-1))$ for $j=1,\dots,40$ and $m=2,\dots,30$. In double precision, $3580209$ triples ($b=2,\dots,10$, $\mathfrak q$ from $0.0005$ to $0.9995$ in steps of $0.0005$, $m=2,\dots,200$) gave no positive $\Delta_{\mathrm H}$ for $\mathfrak q\ge1/(2b-1)$, an upward-closed set of $m$ with $\Delta_{\mathrm H}>0$ for every smaller $\mathfrak q$, and $m^\ast=2$ exactly below $b/(2b^2-1)$; for $b=2$ the zeros of $\Delta_{\mathrm H}(2,\cdot,m)$ are $0.2857$, $0.3062$, $0.3182$, $0.3321$ and $0.3333$ at $m=2$, $3$, $4$, $8$ and $16$. In explicit media with $(b,n_{\mathrm H},A_{\mathrm h})=(2,8,1)$ and $\mathfrak q\in\{0.3,1/3,0.36,0.45\}$, and with $(3,5,1)$ and $\mathfrak q\in\{0.19,0.2,0.25\}$, the contrast of the next-level blocks and the exchangeable baseline, both computed from the iterated coarse matrices at $121$ leaks in $[10^{-4},10^2]$, agreed with the criterion of (b) at every leak and level: at $\mathfrak q=0.3$ and $0.19$, where $m^\ast=3$, the blocks of levels $1$ to $n_{\mathrm H}-2$ qualify above their thresholds and those of level $n_{\mathrm H}-1$ never do, and at the other five values no block of any level qualifies under either judgment.
\end{numcheck}

\begin{numcheck}[Memory of partitioned populations]\label{chk:ren-memory}
Square blocks on tori of $4\times4$ blocks, internal nearest-neighbor coupling $1$, coupling $a_{\mathrm{out}}$ across blocks and uniform leak $\varepsilon^\Pi$, were computed for (side, $a_{\mathrm{out}}$, $\varepsilon^\Pi$) $=(6,0.01,10^{-3})$, $(6,0.001,10^{-3})$, $(4,0.01,10^{-2})$, $(8,0.002,10^{-3})$ and $(6,0.01,0.1)$. The identity of \cref{thm:ren-memory}(a) held to $2.2\times10^{-16}$ and every bound of (a) and (c) held. At $(6,0.01,10^{-3})$, contrast $35.1$, $\Vert\mathcal D\Vert=1.95\times10^{-4}$ against the bound $6.6\times10^{-4}$, $\sin\Theta=0.0244$ against $0.0524$, the relative fast weight $1.5\times10^{-5}$ against $1.5\times10^{-4}$, and the drift error $2.0\times10^{-4}$ against $7.0\times10^{-4}$; the nearest-neighbor coupling of the exact coarse precision was $0.9892$ times the Galerkin coupling, and $0.9989$ times at $a_{\mathrm{out}}=0.001$, contrast $161$. The single-block closure bound of \cref{thm:agy-h1}, on one block with its exterior held fixed, gave $\Xi\le1.9\times10^{-5}$ against the computed $2.4\times10^{-7}$. An independent implementation on $200$ random partitioned populations with two to four blocks, uniform leak and weak cuts found no violation of any bound of (a), (c) or (d), the largest ratio of $\sin\Theta$ to its bound $0.83$, and $\int_0^\infty\Gamma(u)\mathrm du=\mathcal D$ to the quadrature error $2.4\times10^{-8}$ against $\lvert\mathcal D\rvert=2.8\times10^{-4}$. A further run on $4000$ random populations of two to five blocks with one to six members each, with a leak drawn separately for every agent over several orders of magnitude, found no violation of the three bounds of \cref{thm:ren-memory}(d) in their general form, nor of the bounds on $\Vert\mathcal D\Vert$ and $\lvert\mathcal D_{XY}\rvert$ of \cref{thm:ren-closure}(a), beyond relative rounding errors of $3\times10^{-15}$; in the same populations $\lvert\mathcal D\mathsf n\rvert$, with $\mathsf n=(\lvert X\rvert^{1/2})_X$, was as large as $\Vert\mathcal D\Vert\lvert\mathsf n\rvert$ to three digits, so \cref{lem:ren-galerkin}(c) does not extend to nonuniform leaks. The Frobenius quantity $\sigma^2_{\mathrm{cut}}$ is extensive in the number of blocks while the spectral quantities are not: at $(6,0.01,10^{-3})$ on tori of $4\times4$ and of $6\times6$ blocks, $\sigma^2_{\mathrm{cut}}/(\xin^{\min}\xout)=0.776$ and $1.746$, $\Vert\Pi_{\perp\Phi}\Vert^2/(\xin^{\min}\xout)=0.0862$ on both, and $\Vert\mathcal D\Vert/\xout=0.0255$ and $0.0262$; every block satisfied the bound of \cref{lem:ren-galerkin}(b), with $\sigma_X^2=1.00\times10^{-4}$ against $2.67\times10^{-4}$. For two cliques of $n$ agents with internal coupling $1$, joined by one edge of weight $n/2$, at leak $0.01$ and $n=20$, $50$, $100$ and $200$, the contrast was $39.2$, $98.1$, $196.1$ and $392.2$, the exact coarse coupling was $0.513$, $0.505$, $0.503$ and $0.501$ times the Galerkin coupling, $\Vert\mathcal D\Vert/\xout$ ranged from $0.955$ to $0.978$ and $\Vert\Pi_{\perp\Phi}\Vert^2/(\xin\xout)$ from $1.86$ to $1.95$, and $\sigma_X^2=9.5$, $24.5$, $49.5$ and $99.5$ against the bounds $10$, $25$, $50$ and $100$. At $n=20$ the relative fast weight was $1.7\times10^{-4}$ against the bound $2.7\times10^{-3}$ of \cref{thm:ren-memory}(c), and the drift of the slow part differed from $\Pi_\Phi$ by $0.49$ of its norm.
\end{numcheck}

\begin{numcheck}[Per-particle slow modes]\label{chk:ren-perron}
In the block media of \cref{chk:ren-memory}, with the frame of per-block Perron vectors, at contrast $35.1$: $\Vert\mathcal D^{\mathsf P}\Vert=1.59\times10^{-4}$ against the bound $2.45\times10^{-4}$, $\sin\Theta=0.0237$ against $0.0319$, the Galerkin coupling within $2\times10^{-4}$ relative of the exact one, and the slowest retained rate $1.0368\times10^{-3}$ against the exact $10^{-3}$, with the shift bound $2.5\times10^{-4}$. AR0 and AR1, the latter adding white noise of intensity $2\Pi_{\perp\Phi}^\top\Pi_{\perp\perp}^{-2}\Pi_{\perp\Phi}$, had stationary relative entropies $1.40\times10^{-3}$ and $1.36\times10^{-3}$ to the exact law, with largest relative variance errors of $2.3$ percent; at contrast $161$ they were $1.9\times10^{-6}$. The frame of block averages at contrast $35.1$ retained the slowest rate $10^{-3}$ exactly, with relative entropy $1.52\times10^{-3}$ for AR0. An independent implementation found no violation of the bound of \cref{prop:ren-perron}(c) in $200$ instances, with relative shifts of the slowest retained rate up to $0.52$, and the mean coarse rate at most the mean fine rate in $100$ of $100$ instances.
\end{numcheck}

\begin{numcheck}[Without contrast]\label{chk:ren-no-contrast}
For the homogeneous nearest-neighbor precision on $\Z^2$ with coupling $1$ and square blocks of side $b$, the coarse medium has coupling $b$ between adjacent blocks and leak $b^2\varepsilon^\Pi$, hence the squared second-moment length $1/(b\varepsilon^\Pi)$ per axis in block units. It was compared with the squared length of the block averages, computed from \cref{eq:ren-xi2-flow} and, independently, from the real-space second moment of the covariance of the block averages on a torus of $960\times960$ sites, about $30$ correlation lengths at $\varepsilon^\Pi=10^{-3}$; the two computations gave the same ratios to four decimals. The squared length of the coarse medium was $1.9995$, $3.9950$ and $5.9826$ times that of the block averages for $b=2$, $4$ and $6$ at $\varepsilon^\Pi=10^{-3}$, and $3.5726$ times for $b=4$ at $\varepsilon^\Pi=0.1$. On a torus of $240\times240$ sites, under four correlation lengths in half-width at $\varepsilon^\Pi=10^{-3}$, the real-space moment is truncated and the ratios read $2.41$, $4.82$ and $7.21$. On a ring with blocks of $b$ sites at $\varepsilon^\Pi=10^{-3}$, the squared lengths of the coarse medium, $500$, $250$ and $125$, were $2.000$, $3.995$ and $7.958$ times the exact values $250.06$, $62.58$ and $15.71$ for $b=2$, $4$ and $8$.
\end{numcheck}

\begin{numcheck}[Elimination of internal states]\label{chk:ren-adiabatic}
On three random partitioned drifts the identities of \cref{prop:ren-adiabatic}(a) and (b) held to $4.4\times10^{-16}$, $\lvert\Cov(x_\perp,x_{\mathrm{ret}})\rvert\le5.6\times10^{-17}$, and $\Vert M_{\perp\mathrm r}\Vert$ lay between $0.96$ and $1.68$, so the kernel does not vanish. For \cref{obs:ren-adiabatic-rate}, in the first instance the noise intensity of the memoryless reduction was positive semidefinite at every speed factor, and the difference between its drift $-M_{\mathrm{ri}}\mathsf C_{\mathrm{ir}}$ and the quasi-static drift $M_{\mathrm{ri}}M_{\mathrm{ii}}^{-1}M_{\mathrm{ir}}$ fell from $0.378$ to $0.0100$, from $\mathsf s=10$ on in the same proportion as the error. An independent implementation found the identity residual at most $9.8\times10^{-15}$, $\lvert\Cov(x_\perp,x_{\mathrm{ret}})\rvert\le2.2\times10^{-16}$, for a second instance drawn independently the deviations $4.4\times10^{-2}$, $2.0\times10^{-2}$, $8.3\times10^{-3}$, $3.1\times10^{-3}$ and $9.5\times10^{-4}$ at the same speed factors, and, for (c), a nonzero mean $(-0.0139,0.0429)$ of $x_{\mathrm{ret}}(0.05)$ from $x_{\mathrm{ret}}(0)=0$ and $x_{\mathrm{int}}(0)=z$ with $M_{\mathrm{ri}}z\neq0$.
\end{numcheck}

\begin{numcheck}[The error of the memoryless reduction]\label{chk:ren-adiabatic-rate}
Two drifts with three external, two sensory, two active and three internal states, the sparsity of a particular partition and block-diagonal noise were drawn with seeds $20261111$ and $20261112$ and shifted by a multiple of the identity until $M$, $M_{\mathrm{ii}}$ and $\mathsf A_\infty$ had their spectra to the right of $0.5$ (script \texttt{check.py}, part H; all second moments from Lyapunov equations, lags on a grid of spacing $0.005$ over $[0,30]$). At the speed factors $\mathsf s=1$, $3$, $10$, $30$, $100$, $300$, $1000$ and $3000$ the identity of \cref{prop:ren-adiabatic-rate}(a) held to a relative residual of at most $6.4\times10^{-13}$ and the formula for $M_{\perp\mathrm r}$ of (b) to at most $2.6\times10^{-12}$. The identity of \eqref{eq:ren-adiabatic-error}, with both of its integrals evaluated by Simpson quadrature at $\mathsf s=10$ and $u=2$ (inner step $5\times10^{-4}$, outer step $10^{-3}$), held to $2.9\times10^{-10}$ and $1.4\times10^{-10}$ relative in the two instances, and the quadrature of the inner convolution agreed with its value from the propagator of the pair $(x_{\mathrm{ret}},x_\perp)$ to $1.6\times10^{-10}$ and $8.5\times10^{-11}$. In the first instance the relative error $\sup_u\Vert\mathsf E_{\mathrm F}(u)\Vert/\Vert\Sigma_{\mathrm{rr}}\Vert$ was $7.51\times10^{-2}$, $2.97\times10^{-2}$, $9.17\times10^{-3}$, $3.06\times10^{-3}$, $9.18\times10^{-4}$, $3.06\times10^{-4}$, $9.17\times10^{-5}$ and $3.06\times10^{-5}$, so that $\mathsf s$ times the error was $0.0751$, $0.0891$, $0.0917$, $0.0919$, $0.0918$, $0.0918$, $0.0917$ and $0.0917$, against the coefficient $\mathsf c_{\mathrm F}=0.091733$ of \eqref{eq:ren-adiabatic-coefficient}; the supremum over the lags $[0,3]$ of \cref{obs:ren-adiabatic-rate} gave the same values from $\mathsf s=3$ on. The remainder $\mathsf s^2\sup_u\Vert\mathsf E_{\mathrm F}(u)-\mathsf E_1(u)/\mathsf s\Vert/\Vert\Sigma^\infty_{\mathrm{rr}}\Vert$ rose from $0.080$ to $0.193$ and leveled off, as \eqref{eq:ren-adiabatic-rate} requires, and $\Vert\mathsf A_{\mathrm F}-\mathsf A_\infty\Vert$ fell from $0.847$ to $1.03\times10^{-3}$ in proportion to $1/\mathsf s$ from $\mathsf s=100$ on. The bound $\tau_{\mathrm F}\gamma_{\mathrm F}$ of (b) was $2.59$ at $\mathsf s=1$ and $3.24\times10^{-3}$ at $\mathsf s=3000$, a factor $34.5$ and a factor $105.8$ above the computed error. In the second instance $\mathsf s$ times the error was $0.0232$, $0.0467$, $0.0651$, $0.0723$, $0.0752$, $0.0761$, $0.0764$ and $0.0764$ against $\mathsf c_{\mathrm F}=0.076490$, with the remainder between $0.059$ and $0.136$, and the bound of (b) was a factor $96.5$ above the computed error at $\mathsf s=3000$. For part (e) the matrix $\mathsf N_{\mathrm i}$, instead of the intensity, was multiplied by $\mathsf s$. In the two instances $\Vert\mathsf A_{\mathrm w}-\mathsf A_\infty\Vert$ was $1.196$ and $3.115$, the spectrum of $\mathsf A_{\mathrm w}$ lay to the right of $0.71$ and $0.59$, the identity $\mathsf A_{\mathrm w}\Sigma^{\mathrm w}_{\mathrm{rr}}+\Sigma^{\mathrm w}_{\mathrm{rr}}\mathsf A_{\mathrm w}^\top=2\mathsf N_{\mathrm r}\mathsf N_{\mathrm r}^\top$ held to $5.4\times10^{-15}$ relative, and the limit $\sup_u\Vert\mathsf E_{\mathrm w}(u)\Vert/\Vert\Sigma^{\mathrm w}_{\mathrm{rr}}\Vert$ was $0.041706$ and $0.083694$. At $\mathsf s=10$, $100$, $1000$ and $10000$ the relative error was $0.04301$, $0.04176$, $0.04171$ and $0.04171$ in the first instance and $0.08104$, $0.08356$, $0.08368$ and $0.08369$ in the second, so it did not vanish, and $\Vert\mathsf A_{\mathrm F}-\mathsf A_\infty\Vert$ was $1.042$, $1.178$, $1.194$ and $1.196$ in the first and $2.791$, $3.079$, $3.111$ and $3.114$ in the second, so the drift of the reduction did not approach the quasi-static drift. At $\mathsf s=10000$ the products of $\mathsf s$ with $\Vert\mathsf A_{\mathrm F}-\mathsf A_{\mathrm w}\Vert$, with $\Vert\Sigma_{\mathrm{rr}}-\Sigma^{\mathrm w}_{\mathrm{rr}}\Vert/\Vert\Sigma^{\mathrm w}_{\mathrm{rr}}\Vert$ and with $\sup_u\Vert\mathsf E_{\mathrm F}(u)-\mathsf E_{\mathrm w}(u)\Vert/\Vert\Sigma^{\mathrm w}_{\mathrm{rr}}\Vert$ were $3.158$, $0.1052$ and $0.1013$ in the first instance and $4.443$, $0.1565$ and $0.1404$ in the second, each within half a percent of its value at $\mathsf s=1000$, as the rate $1/\mathsf s$ of (e) requires. For a reversible family with a random symmetric positive definite $\Pi$ on seven retained and three internal states (seed $20261113$), the stationary covariance was $\Pi^{-1}$ at every speed factor to $1.3\times10^{-12}$, the formulas of (d) for $\mathsf Z_{\mathrm F}$ and for $\tr\mathsf E_1(u)$ at $u=1$, $3$, $10$ and $30$ held to $3.6\times10^{-15}$ and to a relative $3.0\times10^{-14}$, with every trace positive, $\mathsf s$ times the error rose from $0.0618$ at $\mathsf s=1$ to $0.08861$ at $\mathsf s=3000$ against $\mathsf c_{\mathrm F}=0.088619$, and the bound of (b) was a factor between $25.7$ and $26.8$ above the computed error. For the scalar instance of (d), $M_{\mathrm{rr}}=1$, $M_{\mathrm{ri}}=2$, $M_{\mathrm{ir}}=1$, $M_{\mathrm{ii}}=3$, $\mathsf N=I$, the matrix $\mathsf Z_{\mathrm F}$ was zero and the relative error was $3.42\times10^{-4}$, $3.61\times10^{-6}$ and $3.63\times10^{-8}$ at $\mathsf s=10$, $100$ and $1000$, so that $\mathsf s^2$ times the error tended to $0.0363$; with $M_{\mathrm{ri}}=1$ instead, $\mathsf Z_{\mathrm F}=0.074074$, $\mathsf c_{\mathrm F}=0.040875$, and $\mathsf s$ times the error was $0.040853$ and $0.040873$ at $\mathsf s=100$ and $1000$. The instance of \cref{obs:ren-adiabatic-rate} itself was not recomputed, since its seed is not part of this record.
\end{numcheck}

\begin{numcheck}[Aggregated codes]\label{chk:ren-codes}
With $K=3$, a rank-two code reading $(\mathbf e_1,\mathbf e_2)$, branching four, five levels and Gaussian angles of spread $1.5$, so that $\E\cos\vartheta=e^{-1.125}=0.325$: rotations about $\mathbf e_1$ gave the second singular value $0.171$, $0.287$, $0.575$, $0.385$ and $0.323$ at levels one to five, the first staying $1$; rotations about $\mathbf e_3$ gave equal singular values $0.603$, $0.320$, $0.284$, $0.321$ and $0.326$; frames drawn from the Haar law on $\mathrm{SO}(3)$ gave $(0.382,0.106)$, $(0.272,0.096)$, $(0.138,0.052)$, $(0.074,0.024)$ and $(0.046,0.036)$; and aggregation after true translation kept $(1,1)$.
\end{numcheck}

\begin{numcheck}[Meta-level relays]\label{chk:ren-relay}
Twelve rank-two codes in $\R^4$ that share one direction gave the averaged hop $\bar P$ with eigenvalues $(1,0.4263,0.3206,0.2531)$, hence rank four. The effective rank of $\bar P^n$ at the relative threshold $0.1$ was $4,3,1,1,\dots$ for $n=1,2,3,4,\dots$ up to $n=30$, nonincreasing with the limit $1=\dim V_\cap$, and a relay through one member at a time in cyclic order gave $2,2,2,2,1,\dots$.
\end{numcheck}

\begin{numcheck}[Effective rank along a relay of three codes]\label{chk:ren-rank-nonmonotone}
For the three planes of \cref{ex:ren-rank-nonmonotone} with $\mathsf o=0.05$ the path products after one to six hops in cyclic order had the two largest singular values $(1,1)$, $(1,0.05)$, $(0.05,0.05)$, $(0.05,1.25\times10^{-4})$, $(3.53\times10^{-3},8.85\times10^{-5})$ and $(2.5\times10^{-3},6.25\times10^{-6})$, and the effective ranks at the relative threshold $0.1$ were $2,1,2,1,1,1$ and then $1$ up to twelve hops; with $\mathsf o=0.3$ and the threshold $0.5$ the effective ranks were the same. The common ground had dimension zero in both cases (script \texttt{check\_e.py}; this part is deterministic). In the cyclic continuation the product over one cycle, restricted to $V_3$ and written in the basis $(\mathbf v,\mathbf w)$, agreed with the Jordan block of the proof to $5.6\times10^{-17}$; the logarithms of the two singular values after $3(n+1)$ hops, $n=0,\dots,30$, agreed with \eqref{eq:ren-rank-cyclic} to $1.1\times10^{-13}$; over the hops $6$ to $3000$ the ratio $\sigma_2/\sigma_1$ was at most $0.0499$ times the bound $1/((1-\mathsf o^2)^2n^2)$ for $\mathsf o=0.05$ and $0.2994$ times it for $\mathsf o=0.3$; and the effective rank was one at every hop from $4$ to $93$ in both cases. For $\mathsf o=0.05$ the rates $m^{-1}\ln\sigma_1$ and $m^{-1}\ln\sigma_2$ were $-1.72428$ and $-2.07031$ at $m=30$, $-1.96187$ and $-2.01246$ at $m=300$ and $-1.99286$ and $-1.99946$ at $m=3000$, on either side of $\frac23\ln\mathsf o=-1.99715$; for $\mathsf o=0.3$ they were $-0.79958$ and $-0.80492$ at $m=3000$, against $-0.80265$. For the scope of the example, $4000$ relays of three random planes in $\R^4$ that share a direction (seed $20261131$) were run for nine hops in cyclic order: the effective rank at the thresholds $0.1$, $0.3$, $0.5$, $0.7$ and $0.9$ increased at no hop, and the largest singular value of every path product was one to $7.9\times10^{-15}$. Among $4000$ relays of three random planes without a common direction (seed $20261132$) the effective rank increased at some hop, for at least one of the five thresholds, in $891$. The inequality $\sigma_k(P_VA)\le\sigma_k(A)$ was violated in none of $20000$ random trials in dimensions two to six (seed $20261133$).
\end{numcheck}

\begin{numcheck}[Literal reading across levels]\label{chk:ren-literal}
With literal consensus on cycles at both levels, four groups of five and six groups of three gave hierarchical limit means $0.2993052160$ and $0.2535794074$, equal to the prediction with the prior counted $N$ times to ten digits, against union posterior means $0.765$ and $0.680$. With one complete group of five among cycles, its level-two weight was $0.357$, $0.607$, $0.877$, $0.982$ and $0.9999$ at $n_{\mathrm c,1}=1$, $3$, $6$, $10$ and $20$ rounds at level one, its weight ratio to another group was $(5/3)^{n_{\mathrm c,1}}$, and the top-level mean tended to that group's own literal mean $2/3$. An independent implementation with other data gave $0.1905046646$ and $-0.0787871578$, equal to the prediction, and the information form gave the union posterior means $0.184390$ and $-0.216733$ exactly.
\end{numcheck}

\begin{numcheck}[Finite-round consensus across levels]\label{chk:ren-consensus-levels}
A hierarchy of $60$ agents with $4$, $3$ and $5$ children per unit at levels one to three, lazy Metropolis matrices on cycles with random chords and random representatives (seed $20261121$), with $\rho_{\perp,J}$ at most $0.75$, $0.5$ and $0.7$ at the three levels, was run on scalar readings of a two-dimensional world (script \texttt{check.py}, part F). The protocol, simulated level by level, agreed with the counts \eqref{eq:ren-consensus-counts} to $2.8\times10^{-16}$ relative. With $(30,15,30)$ rounds, $\mathsf d_l=5.36\times10^{-4}$, $6.10\times10^{-5}$ and $9.02\times10^{-5}$, $\mathfrak k_L=6.87\times10^{-4}$, and the largest $\lvert\mathfrak c_{Ij}-1\rvert$ was $3.88\times10^{-4}$, a fraction $0.564$ of the bound; with $(12,8,10)$ rounds, $\mathfrak k_L=0.2283$ and the largest deviation was $0.0971$. In both cases the counts summed to $N$ to $10^{-13}$, the enclosure of the precision in \cref{prop:ren-consensus-levels}(c) held for every child of the top unit, the identity \eqref{eq:ren-consensus-mean} held to $1.3\times10^{-15}$, and the relative entropy to the union posterior was at most $6.8\times10^{-4}$ times the right side of \eqref{eq:ren-consensus-kl}. A hierarchy of $32$ agents with groups of $3$, $5$, $2$, $6$, $4$, $4$, $1$ and $7$ agents, grouped into $2$, $3$ and $3$ and then into one unit (seed $20261122$, rounds $(40,12,12)$), had $\mathfrak k_L=4.24\times10^{-3}$, largest deviation $1.21\times10^{-3}$, and counts that summed to $N$ plus $1.2\times10^{-3}$ to $4.2\times10^{-3}$, not to $N$. The same hierarchy with no round at the first level, rounds $(0,12,12)$, where $\rho_{\perp,J}^0=1$, had $\mathsf d_1=6$ and $\mathfrak k_L=6.0068$, the largest $\lvert\mathfrak c_{Ij}-1\rvert$ equaled $\mathfrak k_L$ to the seven digits printed, the smallest count was zero, the protocol agreed with \eqref{eq:ren-consensus-counts} to $7.5\times10^{-17}$, and the enclosure of (c), with the lower factor zero, held. With $\mathcal W_J=\frac12I+\frac12\one\one^\top/n_J$ on the first hierarchy and the first child as representative, the largest count minus one equaled $\mathfrak k_L$, and the smallest count equaled $\prod_l(1-\rho_l^{n_{\mathrm c,l}})$, to ten digits: $4.25$ and $0.421875$ for $(2,2,2)$ rounds, $0.3073730469$ and $0.8940124512$ for $(5,4,6)$, $14$ and $0.125$ for one round per level, and $9$ and $0$ for $(0,3,2)$ rounds; with $\mathcal W_J=\one\one^\top/n_J$, where $\rho_l=0$, and $(0,2,1)$ rounds they were $3$ and $0$. For a stable two-dimensional world with $60$ agents on a hierarchy of the same shape redrawn with seed $20261123$, with $\rho_{\perp,J}$ at most $0.75$, $0.5$ and $0.827$ at the three levels, where $\bar F(I-\bar P\Lambda^s)$ had spectral radius $0.540$ and $\mathcal S_{\mathrm{cl}}=2.354$, the stationary root-mean-square deviations of the five children of the top unit from the exact coarse filter, computed from a Lyapunov equation for the pair of the world and the deviation, were between $7.3\times10^{-4}$ and $2.3\times10^{-3}$ at $(30,15,30)$ rounds, where $\mathfrak k_L=1.41\times10^{-2}$, and between $0.036$ and $0.107$ at $(12,8,10)$ rounds, where $\mathfrak k_L=0.766$; they were at most $0.020$ times the right side of \eqref{eq:ren-consensus-dynamic-bound}, and between $0.41$ and $0.43$ times the bound obtained from the recursion \eqref{eq:ren-consensus-dynamic} with the exact second moment of its driving term. The average of the five estimates deviated from the exact coarse filter by $7.3\times10^{-5}$ and $1.3\times10^{-2}$ in root mean square, with averaged counts off by up to $3.9\times10^{-4}$ and $6.7\times10^{-2}$, while for a flat consensus of five rounds on the same agents the averaged counts equaled one to $6.7\times10^{-16}$ although single counts were off by up to $30.8$.
\end{numcheck}

\begin{numcheck}[Laplacian renormalization]\label{chk:ren-lrg}
Two $6\times6$ grids with coupling $1$ joined by six edges of coupling $0.01$ have $\lambda_2(L^\Pi)=0.0032336$ and $\lambda_3(L^\Pi)=0.2679$, and the specific heat has maxima at $\tau_{\mathrm L}=1.24$, $8.32$ and $742$; here and below, maxima of the specific heat were located on a grid of $400$ points per decade over $\tau_{\mathrm L}\in[10^{-3},10^6]$ and refined by bounded scalar maximization. The procedure returned the two grids at $\tau_{\mathrm L}=11.2$ and $103$ and one macronode at $928$. At $\varepsilon^\Pi\in\{10^{-4},10^{-2},1\}$ the macronodes for $\varepsilon^\Pi I+L^\Pi$ with the cutoff applied to the shifted spectrum equaled those for $L^\Pi$, while the unshifted cutoff kept $2$, $0$ and $0$ modes; the criterion $r\ge2$ selected the grids and the union, the grids only, and nothing. In \cref{ex:ren-interleaved}, where $\lambda_2(L^\Pi)=2.2205\times10^{-4}$ and $\lambda_3(L^\Pi)=0.26795$, the specific heat has maxima at $\tau_{\mathrm L}=0.0382$, $1.21$ and $1.08\times10^4$. A scan of the procedure over $\tau_{\mathrm L}\in[10^{-3},10^5]$ at $400$ points per decade, with endpoints refined on linear grids of spacing below $10^{-4}$, returned $A$ and $A'$ as macronodes for $\tau_{\mathrm L}$ in $[0.0167,0.2830)$, $[0.3060,0.3871)$ and $[0.7887,0.9998]$ and for no other value, $A\cup A'$ for $1\le\tau_{\mathrm L}<1/\lambda_2=4503.4$, and $B$ for $1/\lambda_3=3.732\le\tau_{\mathrm L}<4503.4$; no scanned value returned $A$, $A'$ and $B$ together. Every computation of this check keeps the modes with eigenvalue strictly below $1/\tau_{\mathrm L}$. The two endpoints at which a threshold is attained were computed separately with the number of kept modes set by index, so that no floating-point comparison at equality enters (script \texttt{ren\_p2\_lrg\_equality.py}, deterministic): at $\tau_{\mathrm L}=1/\lambda_3$ the strict rule keeps two modes and returns $A\cup A'$ and $B$, while keeping the mode at equality gives three macronodes of sizes $1$, $12$ and $35$ and does not return $B$; at $\tau_{\mathrm L}=1/\lambda_2$ the strict rule keeps one mode and returns the whole population. On a homogeneous $16\times16$ torus the specific heat lay between $0.191$ and $0.245$ for $\tau_{\mathrm L}\in[1,10]$, without a peak; the macronodes at $\tau_{\mathrm L}=0.5$, $1$ and $4$ included one cluster of $154$, $235$ and $252$ sites, and the largest macronode contrast at $\varepsilon^\Pi=0.01$ was $2.01$, below the contrast $16.2$ of the whole torus. In the nested medium with $b=4$, $n_{\mathrm H}=3$, $\mathfrak q=0.2$ and $A_{\mathrm h}=1$, with distinct Laplacian eigenvalues $0$, $0.0533$, $0.3067$ and $1.5733$, the specific heat has maxima at $\tau_{\mathrm L}=2.21$, $11.05$ and $53.3$, the procedure returned exactly the level-one blocks at $\tau_{\mathrm L}=2$ and the level-two blocks at $10$, and the contrasts of these blocks were $5.54$ and $6.53$ at $\varepsilon^\Pi=10^{-3}$ and $1.88$ and $1.22$ at $\varepsilon^\Pi=1$. For \cref{conj:ren-lrg}, a clique $C$ of $n$ agents with unit couplings joined by one unit edge to a path of $n_{\mathrm p}$ sites has $\lambda_2(L_C)=n$, $\cut(C)=1$ and $r_\star(C)=n^2$, and the regions obtained from $C$ by removing one member or adding the adjacent site of the path had ratio at most $n+1$ (a site that is not adjacent gives a disconnected region, of ratio $0$), so $C$ is a strict local maximizer of $r_\star$. For $(n,n_{\mathrm p})=(8,64)$, $(16,256)$ and $(24,400)$, with $r_\star(C)/\ln N=15.0$, $45.7$ and $95.2$, the procedure did not return $C$ at any of $3058$, $439$ and $60$ values of $\tau_{\mathrm L}$ in the interval $(1/n,n)$ of \cref{conj:ren-lrg} (logarithmic grids of $3000$, $200$ and $60$ points, and in the first two cases the geometric midpoints of consecutive thresholds $1/\lambda_k$ inside the interval); at least $8$, $21$ and $27$ modes were kept there, and for $n=16$ and $24$ no two members of $C$ were joined at any scanned value. On logarithmic grids over $[10^{-3},10^6]$ the procedure returned $C$ at $31$ of $2000$ values for $n=8$, all in $[9.75,13.31]$, at one of $200$ values for $n=16$, $\tau_{\mathrm L}=45.5$, and at none of $120$ for $n=24$, so that every returned value lies above the interval (script \texttt{ren\_p2\_lrg\_restricted.py}, ties broken at random with seed $7$). A scan on a grid does not exclude windows narrower than its spacing.
\end{numcheck}

\begin{numcheck}[Coarse contrast]\label{chk:ren-closure}
In $200$ random populations with four to seven equal blocks of two to five agents, couplings drawn within blocks with probability $0.9$ and weights in $[0.5,2]$ and across blocks with probability $0.25$ and weights in $[0,0.2]$, so that the partitions are not equitable, leaks in $[10^{-3},1]$ and a random set $S$ of at least two blocks, $r^{(1)}(S)\ge r(\bigcup S)$ held in every instance, with equality in $13$, and the contrasts computed from $\Pi_\Phi$ and from the coarse medium $\mathsf P\Pi\mathsf P^\top$ agreed to $10^{-9}$ relative.
\end{numcheck}
